\section{The Season}

Advent holds four Sundays, the three Ember Days of December, and the seven days before the Vigil
of the Nativity on which the Office sings the great O antiphons. What follows says what the season's name means and what the Fathers taught under it,
how the Church came to keep the weeks before Christmas, and how the four weeks are laid out in
the books of 1962.

\subsection{\latin{Adventus}: The Three Comings}\label{sec:three-comings}

St Bernard opened his first sermon for the season with the name itself:

\begin{quote}
To-day, my brethren, we celebrate the beginning of the season of advent. The name of this, as of
all our other solemnities, is familiar enough and well known to the entire world. But perhaps
the meaning of the name is not so generally understood. For the unhappy children of Adam, taken
up as they are with the pursuit of things transitory and perishable, pay but little attention
to what is salutary and true.\par
{\smallskip\noindent\small ---Bernard, \textit{In Adventu Domini} I.1, tr. a priest of Mount
Melleray, \textit{St Bernard's Sermons for the Seasons}, vol.~I (Dublin, [1921]).\par}
\end{quote}

The Latin \latin{adventus} (gloss: a coming, an arrival) is the word Bernard and the medieval
commentators use for the Lord's coming, and they use it in the plural. In his fifth sermon
Bernard recalls the two comings he has already preached, \latin{duos significantes adventus},
and then names a third: \latin{Tertius enim quidam adventus est medius inter illos} (gloss: for
there is a certain third coming, between those two).

\subsubsection{Two comings}

The oldest form of the teaching is that Christ comes twice. Cyril of Jerusalem opened the
fifteenth of his catechetical lectures with it:

\begin{quote}
We preach not one advent only of Christ, but a second also, far more glorious than the former.
For the former gave a view of His patience; but the latter brings with it the crown of a divine
kingdom. For all things, for the most part, are twofold in our Lord Jesus Christ: a twofold
generation; one, of God, before the ages; and one, of a Virgin, at the close of the ages: His
descents twofold; one, the unobserved, like rain on a fleece; and a second His open coming,
which is to be. In His former advent, He was wrapped in swaddling clothes in the manger; in His
second, He covereth Himself with light as with a garment. In His first coming, He endured the
Cross, despising shame; in His second, He comes attended by a host of Angels, receiving glory. We
rest not then upon His first advent only, but look also for His second.\par
{\smallskip\noindent\small ---Cyril of Jerusalem, \textit{Catechesis} XV.1, tr. E.~H. Gifford,
NPNF, 2nd ser., vol.~7 (New York, 1894).\par}
\end{quote}

Cyril goes on to find both comings in a single prophet. Malachi, he says, speaks of the first when he writes ``And the Lord whom ye seek shall suddenly come to His
temple; behold one coming,'' and of the second when he continues, ``And the Messenger of the
covenant whom ye delight in. Behold, He cometh, saith the Lord Almighty. But who shall abide the
day of His coming? or who shall stand when He appeareth?'' (\textit{Catechesis} XV.2, tr.
Gifford).

Bernard keeps the same pair before his monks and asks them to hold both in mind at once:

\begin{quote}
It is fitting, my brethren, that we should celebrate this season of advent with all possible
devotion, rejoicing in so great a consolation, marvelling at so great a condescension, inflamed
with love by so great a manifestation of charity. But let us not think of that advent only
whereby the Son of man has ``come to seek and to save that which was lost,'' but also of that
other by which He will come again and will take us to Himself. Would to God you kept these two
advents constantly in your thoughts, revolving them in assiduous meditation, pondering in your
hearts how much we have received by the first, how much we are promised at the second! Would to
God you were able thus to ``sleep among the midst of lots''! For these two comings of the Lord
are the two arms of the Bridegroom, between which the Bride was sleeping when she said, ``His
Left Hand is under my head, and His Right shall embrace me''; because, as we read elsewhere,
``Length of days is in (His) Right Hand, and in (His) Left Hand riches and glory.''\par
{\smallskip\noindent\small ---Bernard, \textit{In Adventu Domini} IV.1, tr. a priest of Mount
Melleray (Dublin, [1921]).\par}
\end{quote}

\subsubsection{The hidden coming}

Bernard's fifth sermon begins in its Latin by recalling that he had spoken of sleeping between
the two comings without saying where the sleeper rests. His answer is a third coming, between the
first and the last, and it is the one his hearers live in:

\begin{quote}
This intermediate place is in fact another advent wherein such as have knowledge thereof enjoy
a most pleasant repose. For whereas the other two are known to all, this is secret. In the first
``He (Christ) was seen on earth and conversed with men,'' when, as He Himself bears witness,
``they both saw and hated'' Him. In the last ``all flesh shall see the Salvation of God,'' and
``they shall look on Him Whom they pierced.'' But in that which intervenes He is hidden, visible
only to the elect who see Him in themselves, and so their souls are saved. Therefore, in the
first He came in the flesh and in weakness, in the second He comes in spirit and in virtue, and
in the third He shall come in glory and majesty.\par
{\smallskip\noindent\small ---Bernard, \textit{In Adventu Domini} V.1, tr. a priest of Mount
Melleray (Dublin, [1921]).\par}
\end{quote}

In this passage ``the second'' is the middle coming, the hidden one, and ``the third'' is the
coming in glory: Bernard counts in order of time, so the coming in spirit stands second. His
Latin goes on to call this middle coming \latin{via quædam} (gloss: a kind of road) by which one
passes from the first coming to the last, and to say that in it \latin{requies est et consolatio
nostra} (gloss: he is our rest and consolation). Section~\ref{sec:fathers} gives the same
sermon's answer to how this coming is received.

Bernard has a second set of three comings in another sermon, and it is a different set. In the
third sermon of Advent he writes: ``There are three distinct comings of the Lord of which I
know, His coming to men, His coming into men, and His coming against men. He comes to all men
indifferently, but comes not into all or against all'' (\textit{In Adventu Domini} III.4, tr.
Mount Melleray). The coming \emph{into} men is the indwelling Christ promised, and Bernard cites
the promise: ``If any man love Me, he will keep My word, and My Father will love him, and We
will come to him and will make Our abode with him.'' The coming \emph{against} men is the
judgment of those who refuse him. The fifth sermon arranges the comings by time; the third
arranges them by what each does to the one who receives it. Both teach that Christ comes now,
inwardly.

\subsubsection{The commentators}

The medieval writers on the liturgy gave this teaching as the reason for the season, and they
counted the comings differently because they did not ask the same question. John Beleth asks
what the coming of Christ is, and answers that it is threefold: in the flesh; in the mind, a
coming that happens every day in human souls through the Holy Spirit; and in majesty, on the day
of judgment (\textit{Summa de ecclesiasticis officiis} LXIV). His daily coming in the mind
answers to the coming Bernard calls hidden. Sicard of Cremona and William Durandus ask what the
Church celebrates in these weeks, and both answer with two. For Sicard the Lord's coming is
celebrated for a twofold cause, his coming in humility and his coming in majesty
(\textit{Mitrale} V.1); Durandus counts four comings of the Lord and says that the Church
celebrates two of them, the coming in the flesh and the coming to judgment (\textit{Rationale}
VI.2.3).

\subsubsection{What the season asks}

The prayers of the season ask the Lord to come. The Collect of the First Sunday begins
\latin{Excita, quaesumus, Domine, potentiam tuam, et veni} (gloss: stir up thy power, O Lord,
and come); from 17 December each of the O antiphons at Vespers asks him to come, \latin{veni};
and the Collect of the Vigil names him both as the Redeemer received with joy and as the Judge
who will come (Section~\ref{sec:last-days}). Bernard's counsel in the fourth sermon sets the
reader's work for the four weeks: to ponder how much we have received by the first coming and how
much we are promised at the second. His fifth sermon adds the work between them, keeping the word
by which the hidden coming is received (Section~\ref{sec:fathers}).
